\documentclass[10pt,a4paper]{amsart}
\usepackage{amsmath,amssymb,amsthm,graphicx}
% For commutative diagrams we use Xy-pic
\usepackage[all]{xy}

% We use 2cell for 2-commutative diagrams.
\xyoption{2cell}
% Theorem environments.
%
\theoremstyle{plain}
\newtheorem{theorem}[subsection]{Theorem}
\newtheorem{proposition}[subsection]{Proposition}
\newtheorem{lemma}[subsection]{Lemma}

\theoremstyle{definition}
\newtheorem{definition}[subsection]{Definition}
\newtheorem{example}[subsection]{Example}
\newtheorem{exercise}[subsection]{Exercise}
\newtheorem{situation}[subsection]{Situation}

\theoremstyle{remark}
\newtheorem{remark}[subsection]{Remark}
\newtheorem{remarks}[subsection]{Remarks}

\numberwithin{equation}{subsection}

% Macros
%
\def\lim{\mathop{\mathrm{lim}}\nolimits}
\def\colim{\mathop{\mathrm{colim}}\nolimits}
\def\Spec{\mathop{\mathrm{Spec}}}
\def\Hom{\mathop{\mathrm{Hom}}\nolimits}
\def\Ext{\mathop{\mathrm{Ext}}\nolimits}
\def\SheafHom{\mathop{\mathcal{H}\!\mathit{om}}\nolimits}
\def\SheafExt{\mathop{\mathcal{E}\!\mathit{xt}}\nolimits}
\def\Sch{\mathit{Sch}}
\def\Mor{\mathop{\mathrm{Mor}}\nolimits}
\def\Ob{\mathop{\mathrm{Ob}}\nolimits}
\def\Sh{\mathop{\mathit{Sh}}\nolimits}
\def\NL{\mathop{N\!L}\nolimits}
\def\CH{\mathop{\mathrm{CH}}\nolimits}
\def\proetale{{pro\text{-}\acute{e}tale}}
\def\etale{{\acute{e}tale}}
\def\QCoh{\mathit{QCoh}}
\def\Ker{\mathop{\mathrm{Ker}}}
\def\Im{\mathop{\mathrm{Im}}}
\def\Coker{\mathop{\mathrm{Coker}}}
\def\Coim{\mathop{\mathrm{Coim}}}

% Boxtimes
%
\DeclareMathSymbol{\boxtimes}{\mathbin}{AMSa}{"02}

%
% Macros for moduli stacks/spaces
%
\def\QCohstack{\mathcal{QC}\!\mathit{oh}}
\def\Cohstack{\mathcal{C}\!\mathit{oh}}
\def\Spacesstack{\mathcal{S}\!\mathit{paces}}
\def\Quotfunctor{\mathrm{Quot}}
\def\Hilbfunctor{\mathrm{Hilb}}
\def\Curvesstack{\mathcal{C}\!\mathit{urves}}
\def\Polarizedstack{\mathcal{P}\!\mathit{olarized}}
\def\Complexesstack{\mathcal{C}\!\mathit{omplexes}}
% \Pic is the operator that assigns to X its picard group, usage \Pic(X)
% \Picardstack_{X/B} denotes the Picard stack of X over B
% \Picardfunctor_{X/B} denotes the Picard functor of X over B
\def\Pic{\mathop{\mathrm{Pic}}\nolimits}
\def\Picardstack{\mathcal{P}\!\mathit{ic}}
\def\Picardfunctor{\mathrm{Pic}}
\def\Deformationcategory{\mathcal{D}\!\mathit{ef}}
\usepackage{cleveref}
\crefname{originallemma}{original lemma}{original lemmas}
\Crefname{originallemma}{Original lemma}{Original lemmas}
\makeatletter
\newlabel{algebra-lemma-support-point}{{00L5}{source}{}{sourcepointer-1}{}}
\newlabel{algebra-lemma-support-point@cref}{{[originallemma][0][]00L5}{[1][0][]source}{}{}{}}
\makeatother
\title{A two-periodic complex formed from two injective endomorphisms with commuting images over a one dimensional Noetherian local ring}
\author{The Stacks Project Authors}
\date{stacks-project:chow.tex:local-snapshot-20261012; GNU FDL 1.2 or later}
\begin{document}
\maketitle
\noindent\textit{Editorial note.} All delivered mathematical text is the author's own, copied from the pinned Stacks Project chapter file chow.tex; no statement or proof step was rewritten, reordered or removed.\par\medskip
\noindent\textit{Editorial note.} Preamble provenance: the pinned chapter file opens with an input of the shared preamble and carries no preamble text of its own, so the project's single Stacks-wide preamble staged at sources/stacks-chow/preamble.tex is excerpted verbatim as two spans, lines 14-18 (Xy-pic setup) and lines 154-217 (the native theorem declarations and mathematical macros). Nothing was added, removed or redefined; the wrapper class and the amsmath/amssymb/amsthm/graphicx base packages are supplied by the assembler.\par\medskip
\noindent\textit{Editorial note.} The counted statement and proof contain no cross reference, no citation and no included file, so the delivered body is self contained.\par\medskip
\noindent\textit{Editorial note.} Canonical location of the counted proposition: Stacks tag 02PX (sources/stacks-chow/pinned-original-source.tex line 17371).\par\medskip
\noindent\textit{Allowed standard prerequisites (editorial).}\begin{itemize}
\item Support of a finite module over a Noetherian local ring and finiteness of length at a minimal prime of the support. Algebra, Lemma labelled algebra-lemma-support-point (Stacks tag 00L5), invoked at proof line 17385. Supplies that M\_q has finite length at a minimal prime q of the support and that a module whose support is the closed point has finite length; the exactness computation is inside the counted span.
\item Two-periodic complexes and their exactness. Same chapter, chow.tex, the definition of a (2, 1)-periodic complex and of exactness used in the statement. Fixes the notion whose two kernel conditions are verified in the counted proof.
\end{itemize}\medskip
\ifdefined\hypertarget\hypertarget{sourcepointer-1}{}\fi
\noindent\textit{Source reference.} Original lemma 00L5: sources/stacks-chow/tags.txt, source lines 775--775. Cross-chapter standard prerequisite located in the project tag list staged in the local snapshot; invoked at proof line 17385 of the counted proof. It is retained only as a navigation pointer; every load-bearing step is inside the delivered span.\par\medskip
\begin{lemma}
\label{lemma-pre-symbol}
Let $A$ be a Noetherian local ring.
Let $M$ be a finite $A$-module of dimension $1$.
Assume $\varphi, \psi : M \to M$ are two injective
$A$-module maps, and assume $\varphi(\psi(M)) = \psi(\varphi(M))$,
for example if $\varphi$ and $\psi$ commute.
Then $\text{length}_A(M/\varphi\psi M) < \infty$
and $(M/\varphi\psi M, \varphi, \psi)$ is an exact
$(2, 1)$-periodic complex.
\end{lemma}

\begin{proof}
Let $\mathfrak q$ be a minimal prime of the support of $M$.
Then $M_{\mathfrak q}$ is a finite length $A_{\mathfrak q}$-module,
see Algebra, Lemma \ref{algebra-lemma-support-point}.
Hence both $\varphi$ and $\psi$
induce isomorphisms $M_{\mathfrak q} \to M_{\mathfrak q}$.
Thus the support of $M/\varphi\psi M$ is $\{\mathfrak m_A\}$
and hence it has finite length (see lemma cited above).
Finally, the kernel of $\varphi$ on $M/\varphi\psi M$
is clearly $\psi M/\varphi\psi M$, and hence the kernel
of $\varphi$ is the image of $\psi$ on $M/\varphi\psi M$.
Similarly the other way since $M/\varphi\psi M = M/\psi\varphi M$
by assumption.
\end{proof}

\end{document}
